\begin{abstracten}
\textbf{Context:} Making design decisions in software architecture is crucial to meeting functional and quality requirements. However, it can be challenging due to uncertainty, implicit knowledge, and integration between various phases of the Software Development Life Cycle. This research explores integrating generative AI tools, specifically GPT models, to aid architectural decision-making by creating and using new prompt patterns through a prompt pattern sequence approach. \textbf{Objective:} This study aims to evaluate the use of generative artificial intelligence tools to assist architects in making software architecture decisions using a sequence of prompt patterns. The objective is to understand the strengths and limitations of these tools to guide their application in this context. \textbf{Methodology:} The methodology includes a complete literature review, a preliminary study, the proposal of five new Prompt Patterns, the development of the Prompt Pattern Sequence approach, and industrial research applying the Prompt Pattern Sequence. The literature review provided valuable insights into design pattern selection and revealed a gap in the absence of generative AIs. The preliminary study focused on using generative AI tools to automatically select design patterns in a specific context and evaluate their accuracy and potential to guide developers effectively. New Prompt Patterns specialized in assisting with architectural decisions have been created, and a Sequence of Prompt Patterns has been proposed to evaluate architectural decisions. The industry study involved projects in which software architects participated to assess the usefulness of information generated by sequencing immediate patterns through generative AI tools through interviews and qualitative analysis. \textbf{Outcomes:} The preliminary study of a generative AI tool to help developers choose design patterns took 56 real-world questions, entered them into ChatGPT, and found that 93\% received accurate answers that shows the tool's potential to guide developers in making informed design decisions. The immediate pattern sequence offers relevant and accurate guidance, and industry studies have revealed the practical benefits and implementation of this approach in real-world contexts, offering valuable insights into the effectiveness, limitations, and perceived usefulness of AI-assisted decision support in software architecture.
\end{abstracten}
